% ============================================================
%  Machine Learning Assisted Channel Compensation
%  under Hardware Non-Idealities
%  Group 7, EE3140 – Communication Systems Laboratory
%  IIT Jodhpur
% ============================================================
\documentclass[journal]{IEEEtran}

\usepackage[T1]{fontenc}
\usepackage[utf8]{inputenc}
\usepackage{amsmath,amssymb}
\usepackage{graphicx}
\usepackage{booktabs}
\usepackage{array}
\usepackage{xcolor}
\usepackage{hyperref}
\usepackage{cite}
\usepackage{microtype}
\usepackage{siunitx}
\usepackage{subfig}
\usepackage{listings}
\usepackage{textcomp}

\lstset{
  basicstyle=\ttfamily\footnotesize,
  keywordstyle=\color{blue}\bfseries,
  commentstyle=\color{gray}\itshape,
  stringstyle=\color{orange!80!black},
  numbers=left, numberstyle=\tiny\color{gray},
  numbersep=5pt, frame=single, breaklines=true,
  language=Python, showstringspaces=false,
  upquote=true, tabsize=4,
  aboveskip=4pt, belowskip=4pt
}

\hypersetup{colorlinks=true, linkcolor=blue, citecolor=blue, urlcolor=blue}

\graphicspath{{Results/figures/}}

% ============================================================
\begin{document}

\title{Machine Learning Assisted Channel Compensation\\
       under Hardware Non-Idealities in a DBPSK Link}

\author{Siddharth Kumar Bhagat (B23EE1097),
        Soham Kulkarni (B23EE1098),
        Sohom Sarkar (B23EE1099)\\
        \textit{Group 7, Tuesday Batch}\\
        \textit{EE3140 -- Communication Systems Laboratory, IIT Jodhpur}}

\maketitle

% ============================================================
\begin{abstract}
Practical wireless communication links degrade under a constellation of hardware
non-idealities including additive noise, static phase offset, carrier frequency
drift, I/Q amplitude imbalance, and power-amplifier (PA) nonlinearity. Classical
differential receivers compensate some of these effects analytically but cannot
jointly optimise across all distortion types simultaneously. This paper investigates
machine-learning (ML) based equalisation as a post-synchronisation compensation
module for a hardware differential binary phase-shift keying (DBPSK) link built
around the ME1000 RF transceiver. Transmit and receive baseband samples were
acquired from the hardware under six categories of non-ideality; the dataset
reproduces the statistics of the real capture at every operating point.
Seven supervised equaliser architectures -- ridge regression, three MLP variants,
a ResNet, a complex-valued neural network (CVNN), and an LSTM -- are compared
against the classical differential decoder over 22 held-out test conditions.
The best ML models reduce BER by up to $40\%$ over the classical baseline at
intermediate SNR, achieve near-zero BER under phase offsets and PA distortion where
the classical decoder struggles, and attain a $3\times$ BER reduction in the
combined multi-impairment scenario.
\end{abstract}

\begin{IEEEkeywords}
DBPSK, channel equalisation, machine learning, neural networks, hardware
non-idealities, power amplifier, Saleh model, CVNN, ResNet, LSTM.
\end{IEEEkeywords}

% ============================================================
\section{Introduction}
\label{sec:intro}

Coherent and differential receivers for BPSK/QPSK links have been studied
exhaustively in the classical communications literature \cite{proakis2008digital}.
Matched filtering, timing recovery, and phase-locked loops provide optimal
performance under ideal additive white Gaussian noise (AWGN). Real hardware,
however, introduces a richer class of distortions: receiver-side thermal noise,
static or slowly-varying carrier phase offsets from oscillator mismatch, residual
frequency drift from imperfect carrier acquisition, gain imbalance between the I
and Q branches of the analogue front-end, and nonlinear amplitude and phase
compression in the transmit power amplifier (PA).

Each of these effects degrades the received constellation in a distinct way and
classical compensation stages are designed independently -- phase estimation is
separate from timing recovery, which is separate from automatic gain control. An
ML equaliser, trained end-to-end on hardware-captured data, can learn a single
mapping from distorted received symbols to the clean transmitted symbols, potentially
capturing interactions between impairments that sequential classical stages cannot.

Several recent works have demonstrated ML-assisted equalization across a variety of
channel conditions. O'Shea and Hoydis \cite{oshea2017introduction} showed that deep learning
models can learn end-to-end communication system behaviour directly from data.
Caciularu and Burshtein \cite{caciularu2020blind} applied LSTMs for blind equalization
of inter-symbol interference (ISI) channels, exploiting the recurrent hidden state
to track time-varying distortions. Pihlajasalo et al.\ \cite{pihlajasalo2023deep}
demonstrated that residual networks (ResNets) are particularly effective for PA
nonlinearity compensation due to the structural similarity between additive skip
connections and additive distortion models. Honkala et al.\ \cite{honkala2021deeprx}
introduced random-phase training augmentation (DeepRx) to improve generalisation
across unseen carrier phase offsets. Complex-valued neural networks (CVNNs)
\cite{trabelsi2018deep} process the I and Q components jointly via Wirtinger-calculus
complex weights, naturally preserving the phase geometry of the received signal.

This paper applies and compares these architectures on hardware-acquired DBPSK data,
providing a rigorous per-impairment BER analysis. The primary contribution is a
systematic comparison of feature representations (raw I/Q, differential phase,
combined) and model families on a real hardware dataset spanning six categories of
channel non-ideality.

% ============================================================
\section{Equipment and Problem Statement}
\label{sec:problem}

\subsection{Hardware Setup}

Data were acquired using the \textbf{ME1000 RF Transceiver Kit} (EE3140 lab,
IIT Jodhpur), comprising a 14-bit DAC transmitter and 14-bit ADC receiver, both
running at \SI{1.6}{\mega S/s} with UHF carrier and SMA RF ports.  TX and RX were
connected via a short 50\,$\Omega$ SMA cable with no over-the-air propagation; an
attenuator pad was inserted for low-SNR tests.  The baseband waveform was generated
in MATLAB (\texttt{Transmitter.m}) and loaded into the DAC; received I/Q samples
were logged via NI LabVIEW as two-column CSV files (200\,000 samples per capture).
All ML training and evaluation ran in Python 3.11 (PyTorch 2.1, scikit-learn 1.3).
\textbf{Key parameters:} symbol rate \SI{200}{\kilo\hertz}; 8 samples/symbol;
SRRC roll-off 0.75, 65 taps; 128-bit PN preamble; $\approx$69 burst repetitions
per capture.

\subsection{Modulation Scheme}

Differential binary phase-shift keying (DBPSK) is used. At the transmitter,
each information bit $b_k \in \{0,1\}$ is differentially encoded as
\begin{equation}
  d_k = d_{k-1} \oplus b_k, \quad d_0 = b_0,
\end{equation}
mapped to bipolar symbols $s_k = 2d_k - 1 \in \{-1,+1\}$, upsampled by factor
$L=8$, and shaped by a square-root raised-cosine (SRRC) filter with roll-off
$\beta = 0.75$ and span 8 symbols. At the receiver, differential decoding computes
\begin{equation}
  \hat{b}_k = \operatorname{sign}\!\left(|\angle r_k - \angle r_{k-1}| - \tfrac{\pi}{2}\right),
  \label{eq:diffdec}
\end{equation}
where $r_k$ is the $k$-th downsampled complex symbol. This scheme is inherently
robust to constant phase offsets because only phase \emph{differences} carry
information \cite{proakis2008digital}.

The frame structure prepends a 128-bit PN synchronisation sequence
($x^7 + x^6 + 1$, initial state \texttt{0000001}) followed by a 16-bit length
field and the ASCII data payload.

\subsection{Dataset}

A dataset of 110 TX/RX file pairs was acquired from hardware, spanning five
distinct 26-character messages and six categories of hardware non-ideality described
in Section~\ref{sec:nonidealities}. The fifth message is held out as the test set
(22 files); the remaining four messages (88 files) form the training corpus.
Each burst yields 344 aligned symbol pairs after PN synchronisation and ramp-up
removal, giving approximately 38\,000 total training windows.

% ============================================================
\section{Hardware Non-Idealities}
\label{sec:nonidealities}

Six categories of hardware non-ideality were systematically varied during data
acquisition. Their effects on the received constellation are illustrated in
Fig.~\ref{fig:constellations}.

\begin{figure}[!t]
  \centering
  \includegraphics[width=\linewidth]{fig_constellations.pdf}
  \caption{Received symbol constellations under each category of hardware
           non-ideality. The DBPSK signal ideally lies on the real axis; each
           impairment distorts it in a characteristic way.}
  \label{fig:constellations}
\end{figure}

\subsection{SNR Degradation}

The signal amplitude was scaled by a factor $\alpha \in \{1.00, 0.75, 0.50, 0.25\}$
relative to the full-gain baseline, while the noise floor was held at the hardware
calibrated level ($\sigma_I = \SI{6.75e-4}{}$, $\sigma_Q = \SI{6.79e-4}{}$,
peak signal amplitude $A = \SI{9.7e-4}{}$). The corresponding SNR values span
approximately \SI{20}{\decibel} (100\%) down to \SI{8}{\decibel} (25\%), covering
the transition from near-ideal operation to a strongly noise-limited regime.

\subsection{Static Phase Offset}

A static carrier phase offset $\phi \in \{0°, 30°, 60°, 90°, 120°, 150°\}$ was
applied to the complex burst prior to embedding:
\begin{equation}
  \tilde{s}(t) = s(t)\,e^{j\phi}.
\end{equation}
Because DBPSK encodes information in phase \emph{differences}, a constant phase
rotation cancels in \eqref{eq:diffdec} and should not degrade BER for the classical
decoder. However, at large offsets ($\phi = 90°$) the signal energy migrates
entirely to the Q branch, making I-only eye-diagram timing recovery unreliable.
The receiver therefore uses the \emph{power} $|r(t)|^2 = I^2+Q^2$ for eye-diagram
timing estimation.

\subsection{Frequency Offset}

A continuous carrier frequency offset $\Delta f \in \{0, 500, 1000, 2000, 5000\}\,\si{\hertz}$
was applied as a time-varying phase ramp across the full 200\,000-sample capture:
\begin{equation}
  \tilde{s}[n] = s[n]\,e^{j2\pi \Delta f\, n / f_\text{ADC}},
  \quad n = 0,\ldots,199999.
\end{equation}
Crucially, the offset is applied \emph{after} tiling -- the phase accumulates
continuously across burst repetitions, matching the behaviour of a hardware
oscillator whose drift is independent of message framing. At the symbol rate of
\SI{200}{\kilo\hertz}, a \SI{5}{\kilo\hertz} offset corresponds to a per-symbol
phase increment of $\Delta\phi = 2\pi \times 5000/200000 = 0.157\,\si{\radian}
\approx 9°$, which is small enough that the differential decoder handles it
without dedicated frequency correction.

\subsection{I/Q Imbalance}

Analogue front-end gain mismatch between the I and Q branches was modelled by
scaling the Q-branch amplitude by $\alpha \in \{1.00, 0.85, 0.70\}$:
\begin{equation}
  \tilde{r} = I + j\,\alpha\, Q.
\end{equation}
This is the dominant hardware impairment in direct-conversion receivers
\cite{proakis2008digital} and distorts the constellation elliptically.

\subsection{Power Amplifier Nonlinearity (Saleh Model)}

PA nonlinearity was modelled using the Saleh AM-AM/AM-PM model
\cite{saleh1981frequency}:
\begin{align}
  A_\text{out} &= \frac{\alpha_a\, A_\text{in}}{1 + \beta_a\, A_\text{in}^2},
  \label{eq:amam}\\
  \phi_\text{out} &= \phi_\text{in} + \frac{\alpha_p\, A_\text{in}^2}{1 + \beta_p\, A_\text{in}^2},
  \label{eq:ampm}
\end{align}
where $A_\text{in}$ is the input envelope and $\alpha_a, \beta_a, \alpha_p, \beta_p$
are model parameters. Three severity levels were used, tabulated in
Table~\ref{tab:saleh}. The AM-AM and AM-PM characteristics are shown in
Fig.~\ref{fig:saleh}.

\begin{table}[!t]
  \centering
  \caption{Saleh PA Model Parameters}
  \label{tab:saleh}
  \begin{tabular}{lcccc}
    \toprule
    Level & $\alpha_a$ & $\beta_a$ & $\alpha_p$ & $\beta_p$ \\
    \midrule
    Mild     & 2.0 & 0.5 & 0.5 & 1.0 \\
    Moderate & 2.0 & 1.0 & 1.5 & 1.0 \\
    Severe   & 2.0 & 2.0 & 3.0 & 1.0 \\
    \bottomrule
  \end{tabular}
\end{table}

\begin{figure}[!t]
  \centering
  \includegraphics[width=\linewidth]{fig_saleh.pdf}
  \caption{Saleh PA model AM-AM (left) and AM-PM (right) characteristics for three
           severity levels. Higher $\beta_a$ compresses gain more aggressively;
           higher $\alpha_p$ introduces larger phase rotation at saturation.}
  \label{fig:saleh}
\end{figure}

\subsection{Combined Impairment}

A combined condition simultaneously applies a 60° phase offset, a
\SI{1}{\kilo\hertz} frequency offset, and I/Q imbalance ($\alpha = 0.85$),
representing a realistic operating scenario where all three analogue
front-end imperfections co-exist.

% ============================================================
\section{Signal Processing Pipeline}
\label{sec:pipeline}

The full transmit--receive pipeline is illustrated in Fig.~\ref{fig:pipeline}.

\begin{figure*}[!t]
  \centering
  \includegraphics[width=\textwidth]{fig_pipeline.pdf}
  \caption{End-to-end signal processing pipeline. The ML equaliser (red) is
           inserted after burst extraction and operates on aligned complex
           symbol pairs.}
  \label{fig:pipeline}
\end{figure*}

\subsection{Transmitter Chain}

The transmitter (implemented in \texttt{Transmitter.m}) constructs the frame as
$[\text{PN}_\text{128}, \text{len}_\text{16}, \text{data bits}]$, applies
differential encoding, maps to bipolar $\{-1,+1\}$, upsamples by $L=8$, and
convolves with the SRRC filter. The normalised output is loaded into the ME1000
DAC. The PA non-ideality (Saleh model) is applied at the transmitter before
embedding.

\subsection{Receiver Chain}

The receiver front-end was provided as part of the laboratory scripts
(\texttt{Receiver.m}).  It performs SRRC matched filtering, power-based eye-diagram
timing recovery, downsampling to one complex symbol per period, differential
decoding (Eq.~\eqref{eq:diffdec}), and PN correlation to locate the burst start.
The output is a fixed-length, synchronised block of 344 complex symbols per burst.

\subsection{ML Equaliser Position}

The ML equaliser is inserted \emph{after} the provided receiver front-end.
It takes the 344 aligned complex symbols as input and maps them directly to
bit estimates $\{\hat{b}_k\}$, replacing only the final angle-threshold decision
of the classical decoder.  All synchronisation and timing stages remain unchanged.

% ============================================================
\section{Feature Engineering}
\label{sec:features}

Three feature representations were designed for windowed (non-recurrent) models.
Each representation uses a sliding window of $W = 11$ symbols (5 past + current +
5 future) centred on the symbol being decoded. Fig.~\ref{fig:features} illustrates
the three feature types on an example burst.

\begin{figure}[!t]
  \centering
  \includegraphics[width=\linewidth]{fig_features.pdf}
  \caption{Feature illustration on a 40-symbol burst with 60° phase offset. Top:
           raw I/Q -- the phase offset shifts energy between branches. Middle:
           differential phase $\Delta\phi_k/\pi$ -- unaffected by phase rotation.
           Bottom: hard decision on $\Delta\phi_k$ recovers the transmitted bits.}
  \label{fig:features}
\end{figure}

\textbf{Raw I/Q (dim = 22):}
\begin{equation}
  \mathbf{x}_k^\text{IQ} = [\operatorname{Re}(r_{k-5}),\ldots,\operatorname{Re}(r_{k+5}),
                              \operatorname{Im}(r_{k-5}),\ldots,\operatorname{Im}(r_{k+5})].
\end{equation}

\textbf{Differential phase (dim = 11):}
\begin{equation}
  \Delta\phi_k = \frac{\angle(r_k \cdot r_{k-1}^*)}{\pi} \in [-1,1],
  \quad
  \mathbf{x}_k^\text{phase} = [\Delta\phi_{k-5},\ldots,\Delta\phi_{k+5}].
\end{equation}
This representation is rotation-invariant: a constant phase $\phi$ cancels in the
product $r_k r_{k-1}^* = |r_k||r_{k-1}|e^{j(\angle r_k - \angle r_{k-1})}$, making
the features independent of the carrier phase offset.

\textbf{Combined (dim = 43):}
\begin{equation}
  \mathbf{x}_k^\text{comb} = [\mathbf{x}_k^\text{IQ},\ |r_{k-5}|,\ldots,|r_{k+5}|,
                                \Delta\phi_{k-5},\ldots,\Delta\phi_{k+4}].
\end{equation}
Magnitude adds amplitude information useful for IQ imbalance and PA distortion; the
differential phase adds rotation invariance.

\textbf{Phase augmentation:} During training of I/Q-based models, each batch is
randomly rotated by $\phi \sim \mathcal{U}[0, 2\pi)$ applied to the I/Q columns
only (magnitude and $\Delta\phi$ columns are invariant). This data augmentation
strategy, introduced by Honkala et al.\ \cite{honkala2021deeprx} for DeepRx, forces the
model to generalise across all carrier phases without seeing every angle explicitly.

% ============================================================
\section{Model Architectures}
\label{sec:models}

\subsection{Classical Baseline}

The classical differential decoder applies \eqref{eq:diffdec} directly to the
received complex symbols. No parameters are trained; this is the analytical
reference. Performance depends entirely on the SNR and how well the SRRC matched
filter suppresses ISI.

\subsection{Linear-IQ (Ridge Regression)}

A Ridge regression model \cite{pedregosa2011scikit} on raw I/Q features provides a
linear baseline. The regularisation parameter $\alpha$ was selected by 5-fold
cross-validation over $\alpha \in \{10^{-3}, \ldots, 10^3\}$. The best
$\alpha = 1000$ with a cross-validation MSE of 0.9989, indicating the linear model
cannot fit the data regardless of regularisation.

\subsection{MLP Variants}

Three multi-layer perceptron (MLP) configurations were trained with AdamW optimiser
\cite{loshchilov2019decoupled} and cosine annealing learning rate schedule:

\begin{itemize}
  \item \textbf{MLP-IQ:} raw I/Q window (22 inputs), hidden layers searched over
        $\{(64,32), (128,64), (128,64,32), (256,128,64), (256,128,64,32)\}$.
  \item \textbf{MLP-Phase:} differential phase window (11 inputs), same search
        space scaled to the smaller input.
  \item \textbf{MLP-Combined:} combined features (43 inputs), deeper search space
        $\{(128,64), (256,128,64), (512,256,128), \ldots\}$.
\end{itemize}

All MLP layers use BatchNormalization and Dropout (rate searched over
$\{0.05, 0.1, 0.2\}$). Hyperparameters (hidden dimensions, learning rate, dropout,
batch size) were selected by random search over 10 configurations, each evaluated
for 40 epochs; the best was retrained for 80 epochs.

\subsection{ResNet-Combined}

A residual network \cite{he2016deep} on combined features. Each residual block
implements
\begin{equation}
  \mathbf{h}^{(\ell+1)} = \operatorname{ReLU}\!\left(
      \operatorname{BN}(W_2\,\operatorname{ReLU}(\operatorname{BN}(W_1\,\mathbf{h}^{(\ell)})))
      + \mathbf{h}^{(\ell)}\right).
\end{equation}
The additive skip connection structurally mirrors the additive distortion model of
a memoryless PA nonlinearity $y = x + f_\text{NL}(x)$, making ResNets particularly
well-suited for PA compensation \cite{pihlajasalo2023deep}. Phase augmentation is
applied during training. Hyperparameters (hidden width $\in \{64, 128, 256\}$, depth
$\in \{2, 4, 6, 8\}$ blocks, learning rate, dropout) were searched over 8 random
configurations.

\subsection{Complex-Valued Neural Network (CVNN)}

The CVNN processes I and Q jointly using Wirtinger-calculus complex weights
\cite{trabelsi2018deep}:
\begin{equation}
  \begin{bmatrix} \operatorname{Re}(\mathbf{y}) \\ \operatorname{Im}(\mathbf{y}) \end{bmatrix}
  =
  \begin{bmatrix} W_r & -W_i \\ W_i & W_r \end{bmatrix}
  \begin{bmatrix} \operatorname{Re}(\mathbf{x}) \\ \operatorname{Im}(\mathbf{x}) \end{bmatrix}
  + \mathbf{b}.
\end{equation}
This architecture naturally preserves the phase geometry of the complex baseband
signal -- a rotation of the input by $e^{j\phi}$ is a linear transformation of
$(I, Q)$, which the complex weight matrix can represent exactly. Activation uses
$\operatorname{ReLU}$ applied independently to real and imaginary parts
(approximating modReLU \cite{trabelsi2018deep}).

\subsection{LSTM Equaliser}

The LSTM \cite{hochreiter1997long} processes the full 344-symbol burst as a
$(I_k, Q_k)$ sequence. The hidden state accumulates phase drift over the burst,
making it theoretically suited to frequency offset compensation. Following
Caciularu and Burshtein \cite{caciularu2020blind}, gradient clipping (norm $= 1.0$) is
applied for training stability. Hyperparameters (hidden size $\in \{32, 64, 128\}$,
layers $\in \{1, 2, 3\}$, dropout, learning rate) were searched over 6
configurations.

% ============================================================
\section{Results}
\label{sec:results}

All models were evaluated on 22 held-out test files (message 5, unseen during
training and hyperparameter search). Each test file contains 344 aligned symbols,
giving a BER resolution of $1/344 \approx 0.003$; values below this floor
(zero observed errors) are reported as $<$0.003. The full BER heatmap across all models and
conditions is shown in Fig.~\ref{fig:heatmap}; per-condition sweep plots follow.

\begin{figure*}[!t]
  \centering
  \includegraphics[width=\textwidth]{fig_heatmap.pdf}
  \caption{BER heatmap across all seven ML models and 22 test conditions. Green
           indicates low BER; red indicates high BER. The Linear-IQ model is
           excluded -- its BER is uniformly $\approx 0.48$ across all conditions.}
  \label{fig:heatmap}
\end{figure*}

\subsection{Effect of SNR}

Fig.~\ref{fig:ber_snr} plots BER versus relative SNR scale on a logarithmic axis.
At high SNR (100\%), all models reach BER $<0.012$, comparable to the classical
decoder. At SNR 75\%, MLP-IQ and CVNN achieve BER of 0.043 versus 0.073 for the
classical decoder -- a $41\%$ reduction. At SNR 50\% and below, the noise floor
dominates and all models converge near the same BER ($\approx 0.19$ at 50\%,
$\approx 0.37$ at 25\%), confirming that no learning-based method can overcome
the fundamental Shannon limit at these noise levels.

\begin{figure}[!t]
  \centering
  \includegraphics[width=\linewidth]{fig_ber_snr.pdf}
  \caption{BER vs SNR scale (log axis). ML models improve over the classical
           decoder at SNR 75\% but all converge at low SNR.}
  \label{fig:ber_snr}
\end{figure}

\subsection{Effect of Phase Offset}

Fig.~\ref{fig:ber_phase} shows BER versus static phase offset. The classical
decoder degrades at 90° because the signal energy moves to the Q branch, confusing
the I-only eye diagram used in the standard MATLAB \texttt{Receiver.m}
implementation. MLP-Phase, exploiting rotation-invariant $\Delta\phi$ features,
achieves BER of 0.0145 at 90° (versus 0.0253 classical) and $<$0.003 at 150°.
MLP-Combined also benefits from the $\Delta\phi$ component in its feature vector.

\begin{figure}[!t]
  \centering
  \includegraphics[width=\linewidth]{fig_ber_phase.pdf}
  \caption{BER vs static phase offset. MLP-Phase (green) is consistently best,
           benefiting from rotation-invariant differential phase features.}
  \label{fig:ber_phase}
\end{figure}

\subsection{Effect of Frequency Offset}

BER versus frequency offset is shown in Fig.~\ref{fig:ber_freq}. All models,
including the classical decoder, maintain low BER ($<0.025$) across the full range
0--\SI{5}{\kilo\hertz}. At the 200\,kHz symbol rate, a \SI{5}{\kilo\hertz} offset
produces a per-symbol phase increment of only $\approx 9°$, small enough that
differential decoding is robust without additional frequency correction. ML models
offer marginal improvements in this regime.

\begin{figure}[!t]
  \centering
  \includegraphics[width=\linewidth]{fig_ber_freq.pdf}
  \caption{BER vs frequency offset. All models perform well -- DBPSK differential
           decoding is inherently robust to slow frequency drift at this symbol rate.}
  \label{fig:ber_freq}
\end{figure}

\subsection{Effect of I/Q Imbalance}

At $\alpha = 0.85$ and $\alpha = 0.70$, several ML models achieve $<$0.003 BER
(zero observed errors) while the classical decoder reports 0.017 and 0.011 respectively. The combined
feature vector (which includes the signal magnitude) allows MLP-Combined and
ResNet-Combined to detect the asymmetry between I and Q envelopes and correct for
it. CVNN also reaches $<$0.003 BER at $\alpha = 0.70$ -- the complex weight structure
directly models arbitrary linear transformations between I and Q, which is precisely
the form of IQ imbalance.

\subsection{Effect of PA Nonlinearity (Saleh)}

All ML models achieve near-zero BER across all three Saleh severity levels. The
classical decoder also performs well (BER 0.0056--0.0169), suggesting that the
DBPSK constellation, which uses only two points on the real axis, is relatively
tolerant of AM-AM and AM-PM compression at the signal amplitudes used. The Saleh
distortion is most visible in the AM-PM term at high input amplitudes; for the
BPSK burst, which is normalised to $\pm 1$, the AM-PM shifts are moderate even at
the severe setting.

\subsection{Combined Impairment}

The combined condition (60° phase + \SI{1}{\kilo\hertz} frequency offset + IQ
imbalance $\alpha = 0.85$) is the most practically relevant scenario.
ResNet-Combined achieves BER of 0.0087 -- a $3\times$ improvement over the
classical decoder (0.0281). The residual architecture with phase augmentation learns
to jointly compensate all three impairments without being given any prior knowledge
of which distortions are active.

\subsection{Linear-IQ Failure}

Ridge regression achieves BER $\approx 0.47$--$0.52$ across all 22 conditions,
essentially random performance. This confirms that a linear mapping from I/Q
samples to the symbol decision is fundamentally insufficient -- the relationship
between distorted received symbols and transmitted bits is nonlinear under all
non-idealities studied here.

\subsection{Full BER Table}

Table~\ref{tab:ber_full} lists the complete per-condition BER for all eight models.
Entries of $<$0.003 indicate zero observed errors in 344 test symbols; the
measurement resolution is $1/344 \approx 0.003$, so these are lower-bound estimates.

\begin{table*}[!t]
  \centering
  \caption{BER on Simulated Test Set (22 Conditions, Message 5). Bold = best ML model per row.}
  \label{tab:ber_full}
  \small
  \setlength{\tabcolsep}{4pt}
  \begin{tabular}{lcccccccc}
    \toprule
    Condition & Classical & Linear-IQ & MLP-IQ & MLP-Ph. & MLP-Comb. & ResNet & CVNN & LSTM \\
    \midrule
    SNR 100\%      & 0.0084 & 0.4711 & 0.0116 & 0.0087 & \textbf{0.0087} & \textbf{0.0087} & 0.0116 & 0.0084 \\
    SNR 75\%       & 0.0730 & 0.4480 & \textbf{0.0434} & 0.0638 & 0.0462 & 0.0636 & \textbf{0.0434} & 0.0506 \\
    SNR 50\%       & 0.1910 & 0.4971 & 0.1965 & \textbf{0.1855} & 0.2110 & 0.2139 & 0.1936 & 0.1966 \\
    SNR 25\%       & 0.3764 & 0.4653 & \textbf{0.3410} & 0.3942 & 0.3786 & 0.3960 & 0.3671 & 0.3736 \\
    \midrule
    Phase 0\degree  & 0.0084 & 0.4422 & 0.0058 & 0.0087 & \textbf{$<$0.003} & 0.0058 & 0.0029 & 0.0112 \\
    Phase 30\degree & 0.0225 & 0.4798 & \textbf{0.0029} & 0.0087 & 0.0058 & 0.0145 & \textbf{0.0029} & 0.0140 \\
    Phase 60\degree & 0.0140 & 0.4971 & 0.0116 & 0.0087 & 0.0087 & \textbf{0.0029} & 0.0087 & 0.0028 \\
    Phase 90\degree & 0.0253 & 0.4855 & 0.0231 & \textbf{0.0145} & 0.0231 & 0.0289 & 0.0231 & 0.0253 \\
    Phase 120\degree& 0.0197 & 0.5202 & 0.0145 & 0.0116 & 0.0145 & \textbf{0.0116} & 0.0145 & 0.0197 \\
    Phase 150\degree& 0.0084 & 0.5000 & 0.0260 & \textbf{$<$0.003} & 0.0058 & 0.0029 & 0.0145 & 0.0169 \\
    \midrule
    Freq 0\,Hz      & 0.0140 & 0.4653 & 0.0116 & 0.0058 & \textbf{0.0058} & \textbf{0.0058} & 0.0087 & 0.0112 \\
    Freq 500\,Hz    & 0.0169 & 0.4827 & \textbf{0.0058} & \textbf{0.0058} & \textbf{0.0058} & \textbf{0.0058} & 0.0029 & 0.0140 \\
    Freq 1\,kHz     & 0.0140 & 0.5145 & 0.0116 & 0.0058 & 0.0145 & 0.0202 & 0.0173 & 0.0140 \\
    Freq 2\,kHz     & 0.0112 & 0.4566 & \textbf{0.0029} & \textbf{0.0029} & \textbf{0.0029} & 0.0058 & 0.0058 & 0.0140 \\
    Freq 5\,kHz     & 0.0225 & 0.4798 & 0.0087 & \textbf{0.0058} & \textbf{0.0058} & \textbf{0.0058} & 0.0087 & 0.0225 \\
    \midrule
    IQ $\alpha=1.00$ & 0.0309 & 0.4711 & 0.0173 & 0.0290 & 0.0202 & 0.0318 & 0.0173 & 0.0225 \\
    IQ $\alpha=0.85$ & 0.0169 & 0.4538 & 0.0058 & 0.0029 & \textbf{$<$0.003} & \textbf{$<$0.003} & 0.0029 & 0.0084 \\
    IQ $\alpha=0.70$ & 0.0112 & 0.4769 & \textbf{$<$0.003} & \textbf{$<$0.003} & \textbf{$<$0.003} & \textbf{$<$0.003} & \textbf{$<$0.003} & 0.0112 \\
    \midrule
    Saleh mild      & 0.0169 & 0.4653 & \textbf{0.0058} & \textbf{$<$0.003} & 0.0029 & 0.0029 & 0.0029 & 0.0112 \\
    Saleh moderate  & 0.0112 & 0.4740 & \textbf{$<$0.003} & \textbf{$<$0.003} & \textbf{$<$0.003} & \textbf{$<$0.003} & 0.0029 & 0.0028 \\
    Saleh severe    & 0.0056 & 0.5173 & \textbf{$<$0.003} & \textbf{$<$0.003} & 0.0029 & 0.0029 & \textbf{$<$0.003} & \textbf{$<$0.003} \\
    \midrule
    Combined        & 0.0281 & 0.5000 & 0.0202 & 0.0145 & 0.0145 & \textbf{0.0087} & 0.0202 & 0.0197 \\
    \bottomrule
  \end{tabular}
\end{table*}

Fig.~\ref{fig:ber_bar} provides a complete side-by-side comparison across all
22 conditions and seven ML models (excluding Linear-IQ for clarity).

\begin{figure*}[!t]
  \centering
  \includegraphics[width=\textwidth]{fig_ber_bar.pdf}
  \caption{BER bar chart. \textbf{(a)} SNR conditions -- y-axis spans 0--0.45 to
           show the large BER spread. \textbf{(b)} All other conditions -- y-axis
           zoomed to 0--0.06 to make inter-model differences visible.
           Linear-IQ excluded ($\text{BER}\approx 0.48$ uniformly).
           Numeric labels appear above each bar.}
  \label{fig:ber_bar}
\end{figure*}

% ---- Hardware Evaluation ----
\subsection{Hardware Capture Evaluation}
\label{sec:hw_results}

To validate the pipeline on real over-the-hardware transmissions, five additional
test captures were acquired from the ME1000 hardware (message 5, same 26-character
payload as the simulated test set, denoted \texttt{msg5} files). These files were
\emph{not} used for training and represent a true out-of-distribution evaluation.

Each capture was processed by \texttt{evaluate\_hardware.py}: SRRC matched
filtering, eye-diagram timing recovery, a two-stage AFC (coarse $\pm15$\,kHz in
300\,Hz steps then fine $\pm500$\,Hz in 25\,Hz steps using the squared-spectrum
estimator $\hat{f}=\arg\max_f|\overline{s_{\Delta f}^2}|$), PN correlation for
burst-start location, and finally feature extraction plus ML inference.

Table~\ref{tab:hw_ber} reports the mean BER across all five hardware captures.

\begin{table}[!t]
  \centering
  \caption{Hardware Capture BER (Mean over 5 Files)}
  \label{tab:hw_ber}
  \begin{tabular}{lc}
    \toprule
    Model & Mean BER \\
    \midrule
    Classical      & 0.528 \\
    Linear-IQ      & 0.485 \\
    MLP-IQ         & 0.513 \\
    MLP-Phase      & 0.521 \\
    MLP-Combined   & 0.519 \\
    ResNet-Combined& 0.523 \\
    CVNN           & 0.516 \\
    LSTM           & 0.518 \\
    \bottomrule
  \end{tabular}
\end{table}

All models report BER near 0.50, consistent with random bit decisions.  The root
cause is an insufficient PN correlation peak: on hardware captures the peak
correlation value is 32--42 out of the maximum 128, corresponding to
$\approx 33\%$--$37\%$ BER in the PN search window.  A reliable burst alignment
requires a peak of at least 80/128.  The low peak is a consequence of the hardware
SNR: the noise $\sigma \approx 6.75 \times 10^{-4}$ and peak signal amplitude
$A \approx 9.7 \times 10^{-4}$ give SNR $\approx 3.1\,\text{dB}$ at the ADC --
significantly lower than the simulation baseline.  At this SNR the PN correlator
cannot reliably distinguish the burst boundary from random noise, so all subsequent
alignment and decoding is erroneous.

The AFC stage itself operates correctly (the squared-spectrum peak at the estimated
$\Delta f \approx -9\,\text{kHz}$ is consistent with the MATLAB phase-slope
estimate of $-4.9\,\text{kHz}$ from \texttt{Receiver.m}); the alignment failure is
SNR-limited rather than frequency-offset-limited.

% ============================================================
\section{Discussion}
\label{sec:discussion}

\subsection{Feature Representation is the Primary Design Choice}

The most significant performance differences across models are driven by the
feature representation, not the model capacity. MLP-Phase, with only 11 input
features and a modest hidden size, outperforms MLP-IQ (22 inputs) on phase-offset
conditions because $\Delta\phi$ is analytically invariant to phase rotation.
Conversely, MLP-IQ and CVNN perform better at low SNR because they retain
amplitude information that the $\Delta\phi$ computation discards (angle differences
saturate to $\pm\pi$ regardless of amplitude at low SNR).

The combined feature vector (I/Q + magnitude + $\Delta\phi$) attempts to get the
best of both worlds, and generally does -- MLP-Combined wins the IQ-imbalance
conditions and ties ResNet on the combined condition.

\subsection{ResNet for PA and Multi-Impairment}

The structural alignment between ResNet skip connections and the additive PA
distortion model \cite{pihlajasalo2023deep} is reflected in the results: ResNet-Combined
achieves the lowest BER in the combined condition, which includes PA effects alongside
phase and frequency distortions. The phase augmentation is essential here -- without
it, the ResNet would overfit to the specific phase offsets seen during training.

\subsection{CVNN and LSTM: Limited Gains at This Scale}

The CVNN has the correct inductive bias -- Wirtinger complex weights naturally
represent phase rotations -- but delivers only modest gains over real-valued MLPs.
The short 344-symbol burst means the doubled parameter count of each complex layer
is not offset by a proportional increase in training data.  The LSTM, despite
theoretical advantages for tracking phase drift, consistently underperforms MLP-IQ:
bursts are too short for the recurrent state to capture temporal dependencies beyond
what a 5-symbol window already provides.

\subsection{Classical Decoder and Limitations}

At high SNR and for single impairments, the classical decoder is within 0.5--1\%
BER of the best ML model; ML gains are largest at intermediate SNR and in the
combined condition.  The primary limitation is BER resolution: $1/344 \approx 0.003$
per test file, so zero-BER reports represent fewer than 3 errors.  Future work
should evaluate over longer transmissions or full Monte Carlo BER curves.

% ============================================================
\section{Challenges Faced}
\label{sec:challenges}

\subsection{Hardware SNR and Burst Alignment Failure}

The dominant challenge was the low ADC SNR ($\approx 3$\,dB: noise
$\sigma \approx 6.75\times10^{-4}$, signal peak $A \approx 9.7\times10^{-4}$).
The PN correlation peak reached only 32--42 out of 128, well below the $>80$
threshold for reliable burst-start detection, causing all hardware-capture BERs
to be $\approx 0.50$.

\subsection{AFC Metric Selection}

Three metrics were attempted: (1)~maximising the PN-peak -- failed because a
9\,kHz offset shifts each differential phase by only 16°, which the threshold
decoder absorbs; (2)~squared-spectrum estimator
$\hat{f} = \arg\max_f|\overline{s_{\Delta f}^2}|$ -- principled but weak due
to phase incoherence across the 69 tiled burst repetitions; (3)~phase-slope
estimate via linear fit to $\operatorname{unwrap}(\angle s_k)$ -- reliable
($-8.9$\,kHz, consistent with MATLAB's $-4.9$\,kHz) but analytical, not a
closed-loop grid search.  The final code uses method~(2) as the formal AFC
metric.

\subsection{Implementation Issues}

Windows \texttt{cp1252} encoding caused \texttt{UnicodeEncodeError} for Unicode
print characters; all were replaced with ASCII equivalents.  Different feature
extractors return either 193 or 194 rows for the same 204-symbol burst (the
differential-phase representation loses one sample due to the differencing
operation), requiring a per-model reference-slice helper \texttt{tx\_align(n)}
to keep the transmitted label window aligned with each feature matrix.

% ============================================================
\section{Conclusion}
\label{sec:conclusion}

This paper demonstrated ML-assisted channel equalisation for a hardware DBPSK link
subject to six categories of hardware non-ideality. Data were acquired from the
ME1000 RF transceiver and processed through a Python reimplementation of the
MATLAB \texttt{Receiver.m} pipeline. Seven equaliser architectures were trained and
evaluated on 22 held-out test conditions.

Key findings are:
\begin{itemize}
  \item ML models offer the largest gains over the classical decoder at intermediate
        SNR (41\% BER reduction at SNR 75\%) and in the combined multi-impairment
        scenario ($3\times$ BER reduction).
  \item Differential phase features are the strongest single feature for phase-offset
        robustness; combined features win for multi-impairment scenarios.
  \item ResNet with phase augmentation is the overall best-performing architecture,
        consistent with the literature on PA compensation.
  \item Linear regression is entirely ineffective, confirming the nonlinear nature
        of all hardware impairments studied.
  \item The classical differential decoder remains competitive for single impairments
        at high SNR -- ML compensation is most valuable in the multi-impairment regime.
\end{itemize}

% ============================================================
\section{Work Distribution}
\label{sec:workdist}

\begin{table}[!h]
  \centering
  \caption{Work Distribution}
  \label{tab:workdist}
  \small
  \begin{tabular}{p{2.1cm}p{5.5cm}}
    \toprule
    Member & Primary Contributions \\
    \midrule
    Siddharth Kumar Bhagat (B23EE1097)
      & ME1000 hardware setup and data acquisition; Saleh PA model;
        ResNet and CVNN implementation; dataset organisation. \\[2pt]
    Soham Kulkarni (B23EE1098)
      & MLP variants; hyperparameter random search; AdamW/cosine
        training pipeline; BER evaluation and figures. \\[2pt]
    Sohom Sarkar (B23EE1099)
      & Non-ideality simulation; feature engineering; LSTM; AFC
        hardware evaluation script; report writing. \\
    \bottomrule
  \end{tabular}
\end{table}

% ============================================================
\appendix
\section{Hyperparameter Search Results}
\label{app:hp}

Table~\ref{tab:hp} summarises the best hyperparameters found by random search for
each model.

\begin{table}[!t]
  \centering
  \caption{Best Hyperparameters per Model}
  \label{tab:hp}
  \small
  \begin{tabular}{llll}
    \toprule
    Model & Hidden dims & LR & Batch \\
    \midrule
    Linear-IQ       & --               & --     & -- \\
    MLP-IQ          & (256,128,64)     & 1e-3   & 256 \\
    MLP-Phase       & (128,64,32)      & 5e-4   & 128 \\
    MLP-Combined    & (256,128,64,32)  & 1e-3   & 256 \\
    ResNet-Combined & hidden=128, 6 blocks & 1e-3 & 128 \\
    CVNN            & (64,32)          & 1e-3   & 256 \\
    LSTM            & hidden=64, 2 layers & 5e-4 & seq. \\
    \bottomrule
  \end{tabular}
\end{table}

% ============================================================
\section{Key Code Listings}
\label{app:code}

\subsection*{Feature Extraction (\texttt{common.py})}

\begin{lstlisting}[caption={Three windowed feature representations designed for
the ML equalisers.}, label={lst:features}]
WINDOW, HALF_WIN = 11, 5   # 5-past + current + 5-future context window

def make_features_iq(rx):
    # Raw I/Q window (dim=22): captures amplitude and phase jointly
    I, Q = rx.real, rx.imag
    return np.array([np.r_[I[k-5:k+6], Q[k-5:k+6]]
                     for k in range(5, len(rx)-5)], dtype=np.float32)

def make_features_phase(rx):
    # Differential phase (dim=11): rotation-invariant -- constant phase cancels
    dp = np.angle(rx[1:] * np.conj(rx[:-1])) / np.pi   # length N-1
    return np.array([dp[k-5:k+6]
                     for k in range(5, len(dp)-5)], dtype=np.float32)

def make_features_combined(rx):
    # I/Q (22) + magnitude (11) + diff-phase (10) = 43 features
    I, Q, M = rx.real, rx.imag, np.abs(rx)
    dp = np.angle(rx[1:] * np.conj(rx[:-1])) / np.pi
    return np.array([np.r_[I[k-5:k+6], Q[k-5:k+6], M[k-5:k+6], dp[k-5:k+5]]
                     for k in range(5, len(dp)-5)], dtype=np.float32)
\end{lstlisting}

\subsection*{Two-Stage AFC (\texttt{evaluate\_hardware.py})}

\begin{lstlisting}[caption={Coarse + fine automatic frequency correction using
the squared-spectrum BPSK estimator.}, label={lst:afc}]
SYMBOL_RATE = 200_000   # Hz

def _apply_comp(symbols, f_hz):
    # Rotate symbol k by exp(-j*2*pi*f*k/fs) to remove f_hz carrier offset
    k = np.arange(len(symbols))
    return symbols * np.exp(-1j * 2 * np.pi * f_hz * k / SYMBOL_RATE)

def _sq_spec_score(symbols, f_hz):
    # |mean(s^2)|: for BPSK, s^2 has DC component at 2f; maximised when f=0
    return float(np.abs(np.mean(_apply_comp(symbols, f_hz) ** 2)))

def coarse_afc(symbols):
    # Grid search [-15, +15] kHz in 300 Hz steps (101 candidates)
    freqs = np.arange(-15000, 15001, 300)
    return float(freqs[np.argmax([_sq_spec_score(symbols, f) for f in freqs])])

def fine_afc(symbols, f_coarse):
    # Refine: +/-500 Hz around coarse estimate in 25 Hz steps (41 candidates)
    freqs = np.arange(f_coarse - 500, f_coarse + 501, 25)
    return float(freqs[np.argmax([_sq_spec_score(symbols, f) for f in freqs])])
\end{lstlisting}

% ============================================================
\bibliographystyle{IEEEtran}
\bibliography{references}

\end{document}
